%%%%%%%%%%%%%%%%%%%%%%%%%%%%%%%%%%%%%%%%
%% MCM/ICM LaTeX Template %%
%% 2026 MCM/ICM           %%
%%%%%%%%%%%%%%%%%%%%%%%%%%%%%%%%%%%%%%%%
\documentclass[12pt]{article}
\usepackage{geometry}
\geometry{left=1in,right=0.75in,top=1in,bottom=1in}

%%%%%%%%%%%%%%%%%%%%%%%%%%%%%%%%%%%%%%%%
% Replace ABCDEF in the next line with your chosen problem
% and replace 1111111 with your Team Control Number
\newcommand{\Problem}{A}
\newcommand{\Team}{2620586}
%%%%%%%%%%%%%%%%%%%%%%%%%%%%%%%%%%%%%%%%

\usepackage{newtxtext}
\usepackage{amsmath,amssymb,amsthm}
\usepackage{newtxmath} % must come after amsXXX

\usepackage[pdftex]{graphicx}
\usepackage{xcolor}
\usepackage{fancyhdr}
\lhead{Team \Team}
\rhead{}
\cfoot{}

\newtheorem{theorem}{Theorem}
\newtheorem{corollary}[theorem]{Corollary}
\newtheorem{lemma}[theorem]{Lemma}
\newtheorem{definition}{Definition}

%%%%%%%%%%%%%%%%%%%%%%%%%%%%%%%%
\begin{document}
\graphicspath{{.}}  % Place your graphic files in the same directory as your main document
\DeclareGraphicsExtensions{.pdf, .jpg, .tif, .png}
\thispagestyle{empty}
\vspace*{-16ex}
\centerline{\begin{tabular}{*3{c}}
	\parbox[t]{0.3\linewidth}{\begin{center}\textbf{Problem Chosen}\\ \Large \textcolor{red}{\Problem}\end{center}}
	& \parbox[t]{0.3\linewidth}{\begin{center}\textbf{2026\\ MCM/ICM\\ Summary Sheet}\end{center}}
	& \parbox[t]{0.3\linewidth}{\begin{center}\textbf{Team Control Number}\\ \Large \textcolor{red}{\Team}\end{center}}	\\
	\hline
\end{tabular}}
%%%%%%%%%%% Begin Summary %%%%%%%%%%%
% Enter your summary here replacing the (red) text
% Replace the text from here ...
\begin{center}
\textcolor{red}{%
Use this template to begin typing the first page (summary page) of your electronic report. This \newline
template uses a 12-point Times New Roman font. Submit your paper as an Adobe PDF \newline
electronic file (e.g. 1111111.pdf), typed in English, with a readable font of at least 12-point type.	\\[2ex]
Do not include the name of your school, advisor, or team members on this or any page.	\\[2ex]
Be sure to change the control number and problem choice above.	\\
You may delete these instructions as you begin to type your report here. 	\\[2ex]
\textbf{Follow us @COMAPMath on X or COMAPCHINAOFFICIAL on Weibo for the \newline
most up to date contest information.}
}
\end{center}
% to here
%%%%%%%%%%% End Summary %%%%%%%%%%%

%%%%%%%%%%%%%%%%%%%%%%%%%%%%%%
\clearpage
\pagestyle{fancy}
% Uncomment the next line to generate a Table of Contents
%\tableofcontents
\newpage
\setcounter{page}{1}
\rhead{Page \thepage\ }
%%%%%%%%%%%%%%%%%%%%%%%%%%%%%%


While some animal species exist outside of the usual male or female sexes, most species are
substantially either male or female. Although many species exhibit a 1:1 sex ratio at birth, other
species deviate from an even sex ratio. This is called adaptive sex ratio variation. For example,
the temperature of the nest incubating eggs of the American alligator influences the sex ratios at
birth.

The role of lampreys is complex. In some lake habitats, they are seen as parasites with a
significant impact on the ecosystem, whereas lampreys are also a food source in some regions of
the world, such as Scandinavia, the Baltics, and for some Indigenous peoples of the Pacific
Northwest in North America.

The sex ratio of sea lampreys can vary based on external circumstances. Sea lampreys become
male or female depending on how quickly they grow during the larval stage. These larval growth
rates are influenced by the availability of food. In environments where food availability is low,
growth rates will be lower, and the percentage of males can reach approximately 78% of the
population. In environments where food is more readily available, the percentage of males has
been observed to be approximately 56% of the population.

We focus on the question of sex ratios and their dependence on local conditions, specifically for
sea lampreys. Sea lampreys live in lake or sea habitats and migrate up rivers to spawn. The task
is to examine the advantages and disadvantages of the ability for a species to alter its sex ratio
depending on resource availability. Your team should develop and examine a model to provide
insights into the resulting interactions in an ecosystem.

Questions to examine include the following:

• What is the impact on the larger ecological system when the population of lampreys can
alter its sex ratio?

• What are the advantages and disadvantages to the population of lampreys?

• What is the impact on the stability of the ecosystem given the changes in the sex ratios of
lampreys?

• Can an ecosystem with variable sex ratios in the lamprey population offer advantages to
others in the ecosystem, such as parasites?

Your PDF solution of no more than 25 total pages should include:
• One-page Summary Sheet.
• Table of Contents.
• Your complete solution.
• References list.
• AI Use Report (If used does not count toward the 25-page limit.)

Note: There is no specific required minimum page length for a complete MCM submission. You
may use up to 25 total pages for all your solution work and any additional information you want
to include (for example: drawings, diagrams, calculations, tables). Partial solutions are accepted.
We permit the careful use of AI such as ChatGPT, although it is not necessary to create a solution
to this problem. If you choose to utilize a generative AI, you must follow the COMAP AI use
policy. This will result in an additional AI use report that you must add to the end of your PDF
solution file and does not count toward the 25 total page limit for your solution.
Glossary
Lampreys: Lampreys (sometimes inaccurately called lamprey eels) are an ancient lineage of
jawless fish of the order Petromyzontiformes. The adult lamprey is characterized by a toothed,
funnel-like sucking mouth. Lampreys live mostly in coastal and fresh waters and are found in
most temperate regions.


Notations
\begin{center}
\begin{tabular}{|c|c|}
\hline
Notation & Description \\
\hline
$L$ & Total number of lampreys \\
$M$ & Number of male lampreys \\
$F$ & Number of female lampreys \\
$r$ & Birth rate of lampreys \\
$d$ & Mortality rate of lampreys \\
$p$ & Probability of a female lamprey giving birth to a male \\
\hline
\end{tabular}
\end{center}


\subsection{Mathematical Model}

The dynamic of the Lamprey population can be described by the following differential equations:

\begin{align*}
\frac{dM}{dt} &= rM(F) - dM \\
\frac{dF}{dt} &= rF(M) - dF \\
\end{align*}

where $rM(F)$ and $rF(M)$ are the birth rates of male and female lampreys, respectively. The
mortality rates $dM$ and $dF$ account for natural deaths and other factors. The probability of a
female lamprey giving birth to a male is represented by $p$.


\section{Problem 1}

Here we want to include the Lotka-Volterra Equation:
\[
\begin{cases}
\frac{dx}{dt} = \alpha x - \beta xy \\
\frac{dy}{dt} = \delta xy - \gamma y
\end{cases}
\]

%%%%%%%%%%%%%%%%%%%%%%%%%%%%%%
\end{document}
\end